\ficha{della Paolera e Taylor (2001), Straining at the Anchor}{dellapaolerataylor2001}

\begin{identificacao}
DELLA PAOLERA, Gerardo; TAYLOR, Alan M. \emph{Straining at the Anchor: The
Argentine Currency Board and the Search for Macroeconomic Stability,
1880-1935}. Chicago: University of Chicago Press, 2001. (NBER Series on
Long-Term Factors in Economic Development).

Monografia de história econômica quantitativa, em inglês, lida em recorte de
seis capítulos, cerca de 130 páginas: introdução (p. 3-36), capítulo 2,
\emph{Anchors Aweigh} (p. 37-64), capítulo 3, \emph{The Baring Crisis}
(p. 67-79), capítulo 5, \emph{Relaunching the Gold Standard} (p. 99-117),
capítulo 6, \emph{Calm Before a Storm} (p. 118-136), e capítulo 8,
\emph{Bailing Out} (p. 165-187). A versão lida é a dos extratos por capítulo
distribuídos pelo NBER, cuja numeração de arquivo (c8834 a c8844) segue a ordem
do volume, não a dos capítulos impressos; a paginação citada é a do livro.
Editora, cidade, ano e ISBN vêm da folha de rosto reproduzida em cada extrato;
a série vem do catálogo, porque não há \emph{front matter} nos arquivos.
\end{identificacao}

\bloco{Problema e tese}
Por que a Argentina, que em 1913 tinha renda por habitante comparável à da
Europa ocidental, escorregou para o subdesenvolvimento relativo? A resposta é
institucional e monetária: a estabilidade da \emph{belle époque} repousou sobre
uma âncora nominal rígida, a Caja de Conversión, mas o desenho financeiro em
que essa âncora foi encaixada era incompleto, porque nada garantia a
conversibilidade dos depósitos. Acertar um elemento da arquitetura não impediu
que falhas em outro derrubassem o conjunto.

\bloco{Argumento}
O livro inteiro é a história de uma busca institucional. Percorre de 1880 a
1935 quatro experimentos monetários argentinos, o bimetalismo de 1881, a Lei
dos Bancos Nacionais Garantidos de 1887, a Caja de Conversión de 1890 e a
conversibilidade de 1899, e mostra que cada um respondeu ao fracasso do
anterior, sob a tensão entre a restrição orçamentária do governo e o trilema de
economia aberta. A Caja cumpriu sua missão por décadas graças a um desenho
robusto; o Banco de la Nación, sob supervisão fraca, ampliou objetivos até
socorrer quase todo o sistema financeiro, e para salvar o banco foi preciso
sacrificar a conversibilidade. É o caso, dizem os autores, em que as más
instituições expulsam as boas.

\begin{enumerate}[leftmargin=*,nosep]
\item A Caja nasceu do colapso, não do projeto. A Lei 2741, de 7 de outubro de
1890, deu-lhe o monopólio de emissão e retirou o privilégio dos bancos, mas não
estabeleceu conversibilidade em ouro. As duas leis não se confundem: a primeira
criou a instituição, a segunda disse como conduzir política monetária e cambial
dentro dela (p. 25, nota 17).
\item A paridade de 1899 foi fixada no câmbio de mercado, não no par legal.
Entre 1891 e 1899 a base ficou fixa e sem lastro, e a deflação gerou juros
reais de 10,4 por cento ao ano em média (p. 31). O debate de 1897 opôs setores
urbanos e comerciais, pela paridade antiga, a exportadores e industriais, pela
acomodação da desvalorização (p. 114). Prevaleceram Silvio Gesell e o
\emph{lobby} exportador, e o novo par ficou em 2,27 pesos papel por peso ouro
(p. 117).
\item O desenho é o de uma caixa ortodoxa. A Lei 3871, de 31 de outubro de
1899, mandou trocar papel por espécie a 44 centavos ouro por peso papel, com
lastro marginal de 100 por cento, oferta de moeda endógena e proibição
explícita de mercado aberto, de redesconto e de socorro de liquidez (p. 120).
Cinco diretores nomeados pelo Executivo com aprovação do Senado e mandato de
cinco anos, e a independência se manteve até agosto de 1914 (p. 120-121). A
Caja começou sem reserva alguma: espécie zero em 1891-99, e a meta de 30
milhões de pesos ouro só foi alcançada em 1910 (p. 121).
\item A separação de funções era o núcleo do arranjo. Conversibilidade externa
era tarefa exclusiva da Caja; conversibilidade interna era domínio do Banco de
la Nación e dos bancos privados, e as duas instituições foram mantidas à
distância de propósito (p. 169). O estatuto do banco do Estado limitava o
crédito ao governo nacional a dois milhões de pesos e proibia usar depósitos
para financiar empréstimos públicos; um decreto de 1892 exigiu reserva mínima
de 25 por cento dos depósitos mais um fundo de garantia de 75 por cento dos
depósitos privados, depositado na própria Caja (p. 108).
\item Na \emph{belle époque} o crédito veio dos bancos, não da Caja. Todo o
aumento da base entre 1900 e 1913 foi lastreado em entrada de ouro (p. 121); a
moeda bancária cresceu 161,5 por cento acumulados, o multiplicador subiu de 1,2
em 1893 para 2,2 em 1912, e a razão de reservas caiu de cerca de 50 por cento
para 30 por cento em 1910 (p. 126-127). Em 1912 a circulação tinha lastro de 63
por cento e o risco país estava em 1,5 ponto percentual (p. 129). Ford é citado
para a leitura de que o padrão-ouro amplificava o ciclo na periferia (p. 129).
\item A crise veio de fora e a suspensão foi lida como contingente. Em 1913 o
Banco da Inglaterra elevou a taxa de desconto de 3,4 para 5 por cento e a
oferta de moeda caiu 4,8 por cento (p. 130). Em agosto de 1914 houve corrida,
feriado bancário, moratória de trinta dias e suspensão dos saques de espécie. A
emissão de emergência foi aprovada contra a posição do Executivo, com piso
legal de 40 por cento de lastro (p. 134). O Banco de la Nación redescontou com
dois terços dos 30 milhões do \emph{fondo de conversión}, e a Caja recusou-se a
usar os poderes de emergência, por julgar que perderia credibilidade (p. 135).
\end{enumerate}

\bloco{Conceitos e definições operacionais}
Caixa de conversão é definida por tarefa, não por estatuto.
\chave{The Conversion Office had a simple and easily monitored task: to
exchange gold in its reserves for paper currency, and vice versa, at a fixed
exchange rate, and to otherwise exercise no independent or autonomous monetary
policy functions}{p. 17}
A Caja é classificada como a primeira caixa de conversão plena em país
independente (p. 17). A distinção decisiva do livro é entre conversibilidade
interna, dos depósitos bancários em moeda corrente, e conversibilidade externa,
da moeda corrente em ouro a um par fixo (p. 21). Lastro marginal integral é o
vínculo mecânico entre reservas e base:
\chave{With such a system of 100 percent marginal gold backing for the
currency, there was a strict and inelastic relationship between variations in
the stock of metallic reserves and variations in the monetary base}{p. 120}
Trilema é a impossibilidade de ter juntos câmbio fixo, mobilidade de capital e
política monetária ativa; suposta a abertura, vira dilema entre os dois últimos
(p. 21).

\bloco{Evidência e método}
Estudo de caso quantitativo, 1880-1935, com séries anuais e, a partir de 1908,
mensais, tiradas das \emph{Memorias de Hacienda}, dos balanços do Banco de la
Nación, de cotações de \emph{The Economist} e do arquivo do Bank of London and
the River Plate. Técnicas: estimativa nova do estoque de espécie em mãos do
público (p. 53), teste da abordagem monetária do balanço de pagamentos
(p. 125), e um VAR mensal de 1908-13 calibrado sobre o modelo de equilíbrios
múltiplos de Dornbusch e Frenkel (p. 182). O contrafactual é o balanço dos
bancos privados sem os redescontos do Banco de la Nación depois de 1914
(p. 172).

\bloco{Agentes e vertentes}
Carlos Pellegrini é o arquiteto político da reforma e o defensor da paridade de
mercado. Vicente Fidel López e Victorino de la Plaza conduzem o \emph{Funding
Loan} de 1891; Juan José Romero recusa a rolagem permanente e negocia o Acordo
Romero com o comitê Rothschild (p. 109). Julio Roca é presidente na sanção da
lei de 1899, e José María Rosa o ministro que a executa. Silvio Gesell antecipa
a armadilha de dívida e deflação e sustenta a nova paridade (p. 115-116). Os
autores se filiam ao trilema de Obstfeld e Taylor e ao padrão-ouro como regra
de compromisso de Bordo e Kydland e Bordo e Rockoff, e tratam Williams,
Prebisch e Ford como clássicos e Díaz Alejandro como matriz do declínio.

\bloco{Críticas e limites}
O livro é sobre a Argentina e não é comparativo. O Brasil aparece uma única
vez, dentro da média de sete países latino-americanos da Figura 1.1 (p. 8), e a
Caixa de Conversão brasileira nunca é mencionada, embora seja contemporânea e
institucionalmente irmã. A tese de que o desenho monetário explica a
estabilidade da \emph{belle époque} convive mal com o reconhecimento dos
próprios autores de que o período foi de bonança externa e expansão do estoque
mundial de ouro (p. 125) e de que o choque de 1914 foi exógeno (p. 136). A
avaliação da Caja é, ainda, retrospectiva a partir do plano de 1991, agenda que
o livro declara desde a introdução.

\begin{dialogo}
\item Procurar se a taxa de 15 pence foi defendida como irrevogável ou como
ponto de partida. Na Argentina a escolha de 2,27 foi justificada como
acomodação ao câmbio de mercado contra a volta ao par legal (p. 117). Testa se
o argumento antideflacionista de tipo Gesell circulou nos diários do Rio e de
São Paulo.
\item Procurar quem atribui à Caixa competência de socorro ao crédito ou de
redesconto, e quem a nega. Os autores tratam a proibição de redesconto e de
emprestador de última instância como definidora (p. 120) e a separação frente
ao banco do Estado como o miolo do arranjo (p. 169). Testa se o debate
brasileiro separava a Caixa do Banco do Brasil e do Tesouro do mesmo modo.
\item Procurar o vocabulário sobre limite de emissão, fundo de resgate e fundo
de garantia. Na Argentina o fundo de garantia era depósito compulsório do banco
do Estado na própria Caja (p. 108), e o \emph{fondo de conversión} virou o
instrumento do socorro de 1914 (p. 135). Testa se os fundos brasileiros
cumpriam função equivalente e se a imprensa percebia que abriam brecha na
regra.
\item Procurar como a suspensão de agosto de 1914 foi noticiada. Na Argentina
os agentes a leram como contingente a eventos internacionais, e o câmbio
depreciou menos de 2 por cento em 1915 (p. 135). Testa se a imprensa brasileira
tratou 1914 como contingência externa ou como fracasso do arranjo, e permite
datar quando o compromisso passou a ser lido como revogável.
\item Procurar menções diretas à Argentina e à Caja de Conversión entre 1906 e
1914. O caso argentino era acompanhado em Londres e comentado em \emph{The
Economist} (p. 28, p. 106). Se os diários brasileiros invocavam o precedente, a
comparação da Parte IV sai das próprias fontes.
\end{dialogo}

\usonocapitulo{Parte IV, como caso comparado principal. Serve a três
afirmações. Primeira, que a Caja de 1890, retomada em 1899, foi a primeira
caixa de conversão plena em país independente (p. 17), e que o Brasil de 1906
opera dezesseis anos depois dela, sob problema análogo de credibilidade
periférica. Segunda, que fixar a paridade no câmbio de mercado, e não no par
legal, foi decisão política argumentada, com clivagem setorial identificável
(p. 114, p. 117), o que dá termo de comparação para os 15 pence contra os 27
pence do par legal brasileiro. Terceira, que o teste do arranjo não é o
desempenho em bonança, e sim o comportamento sob choque:
\chave{having a strong and credible currency board may be no defense against a
crisis if the banking sector is rotten and a nasty shock occurs}{p. 169}
O que o livro não autoriza é a transposição direta. Ele não trata do Brasil,
não menciona a Caixa de Conversão e cita o país uma única vez, dentro de uma
média regional (p. 8). Toda equivalência entre os dois desenhos, taxa, limite
de emissão, fundos, papel do Banco do Brasil e suspensão de 1914, tem de sair
da legislação e das fontes brasileiras, e pode ser de contraste.}
